%% kapa: exam=1 round=20 year=2009 period=1 form=A subject=economics subject_ko=경제원론 seq=10 number=50
%% source: Q-Net 감정평가사 기출문제 (cst00309) artlSeq=5024178 file=extracted/1008/5024178/2009년 제20회 감정평가사 1차시험 기출문제/1교시.hwp
%% conversion: kordoc 4.24.51 -> tools/gen1.py (auto); review=proofread-sonnet; flags=none
\begin{question}{50}
빵과 책만 소비하는 어느 경제에서 가격과 소비량은 다음 표와 같다. 2005년도를 기준년도로 할 때 2009년의 실질GDP와 GDP디플레이터의 값은?

{\small\centering
\begin{tabular}{|c|c|c|c|c|}
\hline
 & 빵 가격 & 빵 소비량 & 책 가격 & 책 소비량 \\
\hline
2005년 & 10 & 6 & 10 & 4 \\
\hline
2009년 & 12 & 4 & 9 & 6 \\
\hline
\end{tabular}\par}

\noindent\hspace*{2em}\makebox[2.6cm][l]{\uline{실질GDP}}\uline{GDP디플레이터}\par
\begin{choices}
  \item \makebox[2.6cm][l]{98}100
  \item \makebox[2.6cm][l]{98}102
  \item \makebox[2.6cm][l]{100}100
  \item \makebox[2.6cm][l]{100}102
  \item \makebox[2.6cm][l]{102}98
\end{choices}
\end{question}
\answer{4}
